\chapter{Trabalhos Relacionados}
\label{cap:trabalhos-relacionados}

A utilização de Modelos de \gls{LLM} na \gls{PCG} em jogos é uma área de
pesquisa em expansão. Diversos trabalhos têm explorado o potencial dos \gls{LLM}s
para criar não apenas elementos de jogo isolados, mas também mundos e narrativas
coesas. Nesta seção, são apresentados trabalhos relevantes que abordam a geração
de conteúdo de jogos com \gls{LLM}s, estabelecendo um paralelo com a proposta deste
TCC e contextualizando suas contribuições. A análise segue uma ordem decrescente
de afinidade, partindo do trabalho mais diretamente relacionado até os que
apresentam semelhanças em aspectos mais específicos da abordagem.

O artigo de \citeonline{nasir2024word2world}, intitulado \textit{Word2World},
apresenta uma abordagem que se assemelha de forma significativa aos objetivos
deste TCC. Os autores propõem um sistema que utiliza \gls{LLM}s para gerar
mundos de jogo 2D jogáveis a partir de uma história inicial. O processo é
realizado em um \textit{pipeline} multifásico que inclui a criação da história,
a extração de informações cruciais (como personagens,
\textit{tiles}\footnote{\textit{tiles} (ou ladrilhos, em português) são pequenas imagens
    ou gráficos que são usados para construir o cenário ou mapa do jogo} e
objetivos), a geração do mundo e um ciclo de \textit{feedback} para avaliação e
refinamento. A principal semelhança com o presente trabalho reside no objetivo
central de conectar uma narrativa de alto nível à geração de um nível de jogo
concreto e coerente, utilizando um \gls{LLM} para mediar ambas as tarefas. A
principal diferença, contudo, está no escopo de aplicação. Enquanto
\textit{Word2World} foca em jogos 2D de visão superior de forma genérica, este
TCC se especializa no gênero \textit{roguelike}, que possui restrições e
características próprias, como a alta rejogabilidade e a estrutura de masmorras,
aspectos que exigem uma adaptação e especialização do processo de geração.

De forma similar, o artigo \textit{MarioGPT}, de
\citeonline{sudhakaran2023mariogpt}, demonstra a viabilidade da geração de níveis
a partir de descrições textuais, focando no gênero de plataforma. O sistema
utiliza um modelo \gls{GPT}-2 ajustado para criar fases do jogo \textit{Super Mario
Bros}\footnote{Site oficial do \textit{Super Mario
Bros} (Acesso em: 23 jan. 2026): \url{https://www.nintendo.com/pt-pt/Jogos/NES/Super-Mario-Bros-803853.html}}
a partir de \textit{prompts} simples, como "muitos canos, poucos inimigos,
elevação baixa". A semelhança com este TCC reside na abordagem
\textit{"text-to-level"}, provando que \gls{LLM}s podem interpretar descrições e
produzir conteúdo estruturado e funcional para um jogo. No entanto, a distinção
fundamental está no gênero e na complexidade estrutural. Enquanto os níveis de
plataforma são tipicamente lineares, os mapas de \textit{roguelike} exigem uma
estrutura de grafo com salas interconectadas, um desafio de conectividade e
\textit{layout} que não é o foco principal de \textit{MarioGPT}.

A dissertação de mestrado de \citeonline{auxtero2023game} propõe o sistema
\gls{GEDC}, explorando a criação de um sistema para o design de ambientes de
jogos, utilizando geração procedural com inteligência artificial para facilitar
a criação de narrativas interativas. O trabalho foca em desenvolver um sistema
intuitivo que permite aos usuários criar e experienciar diferentes histórias,
automatizando a geração inicial do ambiente, mas permitindo customização
posterior. A semelhança com este TCC está na aplicação de \gls{LLM}s para a
geração automática de ambientes de jogo e no objetivo de criar uma ferramenta
acessível para criadores. A diferença reside no foco: o trabalho de
\citeonline{auxtero2023game} é mais abrangente, visando jogos de aventura em
geral e com grande ênfase na interface de customização manual, enquanto este TCC
se concentra nas restrições e particularidades algorítmicas específicas do
gênero \textit{roguelike}.

Por fim, em um artigo publicado na \textit{2024 IEEE Conference on Games} (CoG),
\citeonline{gallotta2024consistent} exploram a geração consistente de
conteúdo através de um \textit{framework} chamado \textit{LLMaker}. A principal
contribuição do trabalho é o uso de \textit{function calling}, no qual o \gls{LLM}
atua como um tradutor de comandos em linguagem natural para chamadas de funções
formais em um sistema de \textit{back-end}. Essa abordagem garante que toda
modificação no conteúdo do jogo (neste caso, \textit{layouts} de masmorras) seja
válida e adira às restrições do domínio, mitigando o problema de alucinações e
inconsistências do \gls{LLM}. A semelhança com este TCC está no reconhecimento
do desafio da consistência na geração de conteúdo e na busca por uma solução que
garanta a validade dos artefatos gerados. O domínio de aplicação, focado em
\textit{dungeon crawlers}, também é bastante próximo ao gênero
\textit{roguelike}. A diferença principal, entretanto, está no paradigma de
interação. \textit{LLMaker} é uma ferramenta de \textit{codesign} iterativo, na qual um
designer humano interage continuamente com o sistema via \textit{chat}, enquanto a
proposta deste TCC foca em um processo mais automatizado de geração inicial a
partir de uma entrada textual.

O Quadro~\ref{qua:comparativo_trabalhos} apresenta uma análise comparativa dos
trabalhos relacionados mais próximos. A comparação é baseada em critérios
relevantes para esta pesquisa, destacando as semelhanças e diferenças entre as
abordagens.

\begin{quadro}[h!]
    \centering
    \caption{Comparativo dos Trabalhos Relacionados}
    \label{qua:comparativo_trabalhos}
    \resizebox{\textwidth}{!}{
        \UFCqua{}{
            \begin{tabular}{|l|c|c|c|c|}
                % \toprule
                \hline
                \textbf{Trabalho}             & \shortstack{Geração de níveis                                        \\a partir de descrições textuais} & \shortstack{Foco em engenharia                                        \\de prompt} & \shortstack{Saída em formato\\estruturado (e.g., JSON)} & \shortstack{Abordagem focada\\no gênero Roguelike} \\
                % \midrule
                \hline
                Word2World \cite{nasir2024word2world}    & x                             & x          &            &            \\
                \hline
                MarioGPT \cite{sudhakaran2023mariogpt} & x                             &            &            &            \\
                \hline
                \gls{GEDC} \cite{auxtero2023game}        & x                             &            & x          &            \\
                \hline
                LLMaker \cite{gallotta2024consistent} & x                             & x          &            &            \\
                \hline
                \textbf{Trabalho Proposto}    & \textbf{x}                    & \textbf{x} & \textbf{x} & \textbf{x} \\
                \hline
                % \bottomrule
            \end{tabular}
        }{
            \Fonte{Elaborado pelo autor}
        }
    }
\end{quadro}